\chapter{Introduction}\label{ch:intro}

\section{Goals and Objectives of the Project}
\guide{List the main \textbf{goal} (some sentences) and 5-8 \textbf{specific, measurable} objectives (SMART).
Example: ``Reduce manual attendance recording time from 10 min to under 1 min per class.''}
\todo{Write goals and objectives.}
\begin{enumerate}
  \item \todo{Objective 1 (measurable)}
  \item \todo{Objective 2 (measurable)}
\end{enumerate}

\section{Scope of the Work}
\guide{State clearly what the system \textbf{will} do (in scope) and what it \textbf{will not} do (out of scope).}
\subsection{Current Situation and Context}
\guide{Describe the existing problem or manual system and why a new system is needed. Use facts, statistics or interviews/surveys with citations.}
\todo{Describe the current situation.}

\subsection{Competing Products (Available in Market)}
\guide{Compare 3--5 existing products. Summarise features and \emph{limitations} in a table, then explain how your system differs.}
\begin{table}[htbp]
\centering
\caption{Comparison with existing products}\label{tab:competitors}
\begin{tabularx}{\textwidth}{@{}l X X X@{}}
\toprule
\textbf{Product} & \textbf{Key Features} & \textbf{Limitations} & \textbf{Our Advantage}\\
\midrule
Product A & \todo{...} & \todo{...} & \todo{...}\\
Product B & \todo{...} & \todo{...} & \todo{...}\\
\bottomrule
\end{tabularx}
\end{table}

\section{System Overview}
\guide{High-level description: what the system does, main features, target users, and (if applicable) the AI capability. You can include one overview figure.}
\todo{Write system overview.}
\figplaceholder{system_overview.png}{System overview}{fig:overview}

\section{Structure of the Document}
\guide{One or two sentences per chapter describing its contents. Update after finishing the report.}
The remainder of this report is organised as follows. \Cref{ch:pm} presents the project management plan.
\Cref{ch:req} specifies the requirements. \Cref{ch:arch} describes the architecture.
\Cref{ch:design} details the design. \Cref{ch:impl} explains the implementation.
\Cref{ch:test} reports testing and quality assurance. \Cref{ch:deploy} covers deployment and go-live, and
\Cref{ch:sustain} presents the sustainability plan.

\section{Terms, Acronyms, and Abbreviations Used}
\guide{Define every technical term, acronym and abbreviation used in the report (alphabetical order).}
\begin{table}[htbp]
\centering
\caption{Terms, acronyms and abbreviations}\label{tab:terms}
\begin{tabularx}{\textwidth}{@{}l X@{}}
\toprule
\textbf{Term} & \textbf{Definition}\\
\midrule
API   & Application Programming Interface\\
CI/CD & Continuous Integration / Continuous Deployment\\
ER    & Entity--Relationship\\
UAT   & User Acceptance Testing\\
VAPT  & Vulnerability Assessment and Penetration Testing\\
\bottomrule
\end{tabularx}
\end{table}
